\PassOptionsToPackage{dvipsnames,svgnames,table}{xcolor}
\documentclass[10pt]{article}
\usepackage[a4paper,margin=22mm]{geometry}
\usepackage[no-math]{fontspec}\setmainfont[SmallCapsFont={Latin Modern Roman Caps}]{Latin Modern Roman}
\usepackage{amsmath,amssymb,amsthm,mathtools,graphicx,xcolor,hyperref,xurl,xspace,ifthen,enumerate,textcomp}
\usepackage[shortlabels]{enumitem}
\setlength{\emergencystretch}{4em}
\newtheorem{theorem}{Theorem}[section]
\newtheorem{theo}[theorem]{Theorem}
\newtheorem{thm}[theorem]{Theorem}
\newtheorem{lemma}[theorem]{Lemma}
\newtheorem{lemm}[theorem]{Lemma}
\newtheorem{prop}[theorem]{Proposition}
\newtheorem{coro}[theorem]{Corollary}
\newtheorem{corollary}[theorem]{Corollary}
\newtheorem{fact}[theorem]{Fact}
\newtheorem{claim}[theorem]{Claim}
\theoremstyle{definition}
\newtheorem{defi}[theorem]{Definition}
\newtheorem{definition}[theorem]{Definition}
\newtheorem{rema}[theorem]{Remark}
\newtheorem{remark}[theorem]{Remark}
\newtheorem{exam}[theorem]{Example}
\newtheorem{example}[theorem]{Example}
\newenvironment{enonce*}[2][]{\par\medskip\noindent\textbf{#2.}\itshape}{\par\medskip}
\providecommand{\eg}{e.g.\ }
\providecommand{\ie}{i.e.\ }
\providecommand{\resp}{resp.\ }
\providecommand{\cf}{cf.\ }
\providecommand{\longthanks}[1]{\par\smallskip\noindent #1\par\smallskip}
\newcommand{\Z}{{\mathbb Z}}
\numberwithin{equation}{section}
\newcommand{\eee}{\mathrm e}
\newcommand{\expref}[2]{\texorpdfstring{\hyperref[#2]{#1~\ref{#2}}}{#1~\ref{#2}}}
\newcommand{\epfmtdoi}[1]{\href{https://doi.org/#1}{doi:#1}}
\newcommand{\epfmt}[2]{#1:\,#2}
\newcommand{\bin}{{\rm bin}}
\newcommand{\cube}{\{0,1\}^n}
%
%
\newcommand{\bP}{\mathbb{P}}
\newcommand{\cD}{\mathcal{D}}
\newcommand{\cF}{\mathcal{F}}
\newcommand{\Fq}{{\mathbb{F}_q}}
\newcommand{\Ftwo}{{\mathbb{F}_2}}
\newcommand{\Fp}{{\mathbb{F}_p}}
\newcommand{\cL}{\mathcal{L}}
\newcommand{\cP}{\mathcal{P}}
\newcommand{\half}{\frac{1}{2}}
\newcommand{\eps}{\epsilon}
\newcommand{\logeps}{\log({1 \over \epsilon})}
\newcommand{\logepsk}{\log({k \over \epsilon})}
\newcommand{\set}[1]{{\left\{#1\right\}}}
\newcommand{\logepsksquare}{\log^2({k \over \epsilon})}
\newcommand{\ra}{\rangle}
\newcommand{\la}{\langle}
\newcommand{\oneV}{\mathbf{1}}
\newcommand{\Exp}{{\mathbb{E}}}
\newcommand{\poly}{{\rm poly}}
\newcommand{\AllPl}{{\mathcal{P}}}
%
\newcommand{\Pl}{{Y}}
\newcommand{\Tr}{{\rm Tr}}
%
\renewcommand{\div}{{\rm div}}
\newcommand{\closure}{{\rm Closure}}
\newcommand{\supp}{{\rm Supp}}
\newcommand{\con}{{\rm Con}}
\begin{document}
\section*{1542. Length lower bound for balanced binary codes obtained by algebraic-geometric and Hadamard concatenation with designated-distance analysis}
Avraham Ben-Aroya and Amnon Ta-Shma. \emph{Constructing Small-Bias Sets from Algebraic-Geometric Codes}. Theory of Computing volume 9 article 5 (2013), official complete journal PDF acquired 2026-09-30; canonical DOI 10.4086/toc.2013.v009a005 verified against official landing and printed PDF header. \url{https://theoryofcomputing.org/articles/v009a005/}. CC BY3.0.
\paragraph{Selected problem and scope (editorial).} Exact original Theorem4.1 only, for binary balanced codes constructed by geometric Goppa evaluation concatenated with Hadamard and analyzed using the stated designated distance. It is not a bound for every code actual minimum distance. Complete original balanced-code terminology, function-field/place/divisor/Riemann-Roch/Goppa/Hadamard setup, Section3 construction context, and all Section4 descent/conorm constructions and deferred proof are retained. Standard weak approximation, splitting/conorm degree identities, Hasse-Weil and the explicitly stated Castelnuovo bound remain named prerequisites. No upper-bound construction, small-bias reformulation or descent helper is counted separately. Official native archivePDF is byte-identical to the corpus published original.
\paragraph{Difficulty (editorial).} After the named function-field, Castelnuovo and Hasse-Weil prerequisites, the selected designated-distance lower bound still requires the new closure-subfield descent: construct the divisor/evaluation set, preserve Riemann-Roch dimension and compare designated relative distance. The full conorm claim and its deferred valuation proof at1037–1143 are retained, followed by the distinct large-field/few-evaluation-point cases and the final genus/degree/dimension bounds at1148–1226. This source-local descent and parameter-obstruction mechanism constitutes one H2 outcome beyond routine substitution into an AG-code distance bound.
\paragraph{Formatting (editorial).} Exact human-authored mathematical text, formulas and proof order are retained in source order. Only standalone typography/front matter, original counter restoration, bibliography presentation and the declared noncore printed reference locators are adapted. Non-rendered author comments are excluded while every percent newline guard is retained. No mathematical argument, correction or reconstruction is generated.
\paragraph{Original source quirks (editorial).} The exact authored descent recipe retains native1067 wording z in F/K and the original valuation/domain/subscript conventions throughout the construction over F prime. Every surrounding triplet, conorm, dimension and designated-distance passage is supplied without editor domain/index substitution, new argument or mathematical-refereeing conclusion.
\clearpage

The important case of \emph{binary} error correcting codes is still
open. In the binary case, the Gilbert-Varshamov bound gives the best
known (explicit or non-explicit) codes to date. For codes with
distance close to half, the bound shows that random linear codes
of length $n=O(k/\eps^2)$ are $[n,k,(1/2-\eps)n]_2$ codes. Finding
an explicit construction that attains this bound is an open problem.

Another closely related question is that of finding an
$[n,k,(1/2-\eps)n]_2$ binary code, in which the relative weight of
every non-zero codeword is in the range $[1/2-\eps,1/2+\eps]$.
Such codes are called \emph{$\eps$-balanced} and they are related to
another kind of combinatorial objects called \emph{$\eps$-biased
sets}. An $\eps$-biased set is a set $S \subseteq \set{0,1}^k$ such
that for every non-empty subset $T \subseteq [k]$, the binary random
variable $\bigoplus_{i \in T} s_i$, where $s$ is sampled uniformly
from $S$, has bias at most $\eps$. It turns out that $\eps$-biased
sets are just $\eps$-balanced codes in a different guise: the
%
%
rows of a matrix whose columns generate an $\eps$-balanced code form an $\eps$-biased set, and vice
versa. In terms of parameters, an $[n,k]_2$ $\eps$-balanced code is
equivalent to an $\eps$-biased set $S \subseteq \set{0,1}^k$ of size
$n$.

The status of $\eps$-balanced codes is similar to that of
$[n,k,(1/2-\eps)n]_2$ codes. In both cases the probabilistic
method gives non-explicit $[n,k]_2$ $\eps$-balanced codes with
$n=O({k / \eps^2})$, whereas the best lower bound is
\[
n=\Omega\left(\frac{k}{\eps^2 \logeps}\right)\,.
\] For a discussion of these
bounds see~\cite[Section 7]{AGHP:92}.
\setcounter{section}{2}
\section{Restating the construction in the terminology of algebraic function fields}
\label{sec:Restating}

Without putting the above construction in the proper context, it may
appear coincidental. We now describe the general framework of
algebraic-geometric codes and explain why the above construction
fits into this framework.

\subsection{Algebraic geometry}\label{subsec:AG-prel}
We recall a few notions from the theory of algebraic function
fields. A detailed exposition of the subject can be found, \eg,
in~\cite{S:93}.

Let $\Fq $ denote the finite field with $q$ elements. The polynomial
ring $\Fq [x]$, where $x$ is transcendental over $\Fq $, is the set
of all \emph{polynomials} in $x$ with coefficients in $\Fq $. The
\emph{rational} function field $\Fq (x)$, where $x$ is
transcendental over $\Fq $, contains all \emph{rational} functions
in $x$ with coefficients in $\Fq $. A field $F$ is an algebraic
function field over $\Fq$, denoted $F/\Fq $, if $F$ is a
\emph{finite} algebraic extension of $\Fq (x)$.

A \emph{place} $P$ of $F/\Fq $ is a maximal ideal of some valuation
ring $O$ of the function field. We denote by $O_P$ the valuation
ring that corresponds to the place $P$. We denote by $v_P$ the
\emph{discrete valuation} that corresponds to the valuation ring
$O_P$. Therefore, we can write $P$ and $O_P$ as $$P=\set{y \in F \;:\;
v_P(y)  >0}\quad\text{and}\quad O_P=\set{y \in F \;:\; v_P(y) \ge
0}\,.$$ Since $P$ is a maximal ideal, $F_P=O_P/P$ is a field. In fact,
it is a finite field~\cite[Proposition I.1.14]{S:93}. For every $y
\in O_P$, $y(P)$ denotes $y \pmod P$ and is an element of $F_P$. It
can be thought of as the evaluation of the function $y$ at the
``evaluation point'' $P$. The degree of a place $P$ is defined to be
$\deg(P)=[F_P:\Fq ]$, \ie, the dimension of $F_P$ as a vector space
over $\Fq$. In particular, if a place $P$ is of degree $1$ then
$F_P$ is isomorphic to $\Fq$ and the evaluation of $y$ at the place
$P$ is an element of $\Fq$.

We proceed with a simple example that illustrates the above notions.
\begin{example}\label{exa:rational}
The rational function field $\Fq(x)/\Fq$ is the simplest algebraic function field.
For every irreducible polynomial
$p(x)$ let $O_p$ denote the set of rational functions
$r(x)={u(x)}/{w(x)}$ whose denominator $w(x)$ is not divisible
by $p$. The set $O_p$ forms a ring. Furthermore, for every rational function
$r$, either $r$ or $r^{-1}$ belongs to $O_p$, making $O_p$ a
valuation ring. The place $P_p$ is the set of rational functions
$r(x)={u(x)}/{w(x)}$ for which $p$ divides $u$ but does not
divide $w$.

There is exactly one more valuation ring in $\Fq(x)$.
Let $O_\infty$ denote the set of rational functions
$r(x)={u(x)}/{w(x)}$ for which $\deg(w) \ge \deg(u)$. Again,
$O_\infty$ is a ring and furthermore, for every rational function
$r(x)= {u(x)}/{w(x)}$, either $r$ or $r^{-1}$ belongs to
$O_\infty$, making $O_\infty$ a valuation ring. The place $P_\infty$
is the set of rational functions $r(x)= {u(x)}/{w(x)}$ for which
$\deg(w) > \deg(u)$.

For an irreducible polynomial $p$, the discrete valuation $v_p$ corresponding to $O_p$ is defined as follows.
For a polynomial $u(x) \in \Fq[x]$, $v_p(u)$ is the largest integer $k$ such that $p^k$ divides
$u$. For a rational function $r(x)= {u(x)}/{w(x)}$,
$v_p(r)=v_p(u)-v_p(w)$. Thus, $v_p(r)$ counts the number of zeroes
(or poles) that $r$ has when the irreducible polynomial $p(x)$ is
zero. If $p$ is linear, \ie, $p(x)=x-\alpha$, $v_p(r)$ counts the
number of zeroes (or poles) the polynomial $p(x)$ has when setting
$x=\alpha$. If in addition $r=u/w$ belongs to $O_p$, \ie, $p$ does
not divide $w$, then $r(P_{x-\alpha})=r \bmod (x-\alpha)$ is well
defined, and, in fact, is the element $u(\alpha)/w(\alpha)$ in
$\Fq$.

The discrete valuation $v_\infty$ corresponding to $O_\infty$ is defined as
follows. For a polynomial $u(x) \in \Fq[x]$, $v_\infty(u)$ is
$-\deg(u)$. For a rational function $r(x)= {u(x)}/{w(x)}$,
$v_\infty(r)=v_\infty(u)-v_\infty(w)=\deg(w)-\deg(u)$. Thus,
$v_\infty(r)$ counts the number of zeroes (or poles) that $r$ has
when $x=\infty$. If $r(x)=u(x)/w(x)$ belongs to $O_\infty$, \ie,
$\deg(w) \ge \deg(u)$, then $r(P_\infty)$ is well defined, and is
either zero (if $\deg(w) > \deg(u)$) or the element $u_k/w_k$ in
$\Fq$, where $u_k$ and $w_k$ are the leading coefficients of the
polynomials $u$ and $w$ respectively.
\end{example}

We let $\AllPl_F$ denote the set of places of $F$, and $N(F)$ the
number of places of \emph{degree $1$} in $F/\Fq $. $N(F)$ is always
finite. $\cD_F$ is the free abelian group over the places of $F$. A
\emph{divisor} is an element in this group, \ie, it is a formal sum
$G=\sum_{P \in \AllPl_F} n_P P$ with $n_P \in \Z$ and where $n_P
\neq 0$ for only a finite number of places. We also denote
$v_P(G)=n_P$. The \emph{degree} of the divisor $\sum_P n_P P$ is
defined to be $\sum_P n_P \cdot \deg(P)$, and is always finite.
We say $G_1 \ge G_2$ if $G_1$ is component-wise larger than $G_2$,
\ie, $v_P(G_1) \ge v_P(G_2)$ for every place $P$. The \emph{support}
of a divisor $G$ is $\supp(G) = \set{ P \in \AllPl_F \;:\; v_P(G)\neq
0}$.

Each element $0 \neq x \in F$ is associated with two divisors. The
first is called the \emph{principal divisor} of $x$ and it is defined
by
$$(x) = \sum_{P} v_P(x) P\,.$$
The degree of a principal devisor is always $0$. The second is the
\emph{pole divisor} of $x$ and it is defined by
$$(x)_\infty = \sum_{P: v_P(x) < 0} -v_P(x) P\,.$$ If $x\in F \setminus \Fq$ then $\deg((x)_\infty)=[F:\Fq (x)]$.

\begin{example}[Continued from \expref{Example}{exa:rational}]
\label{exa:rational:2} Let $u \in \Fq[x]$ be an arbitrary nonconstant polynomial.
Then, $u$ has a pole at $P_\infty$ (since
$v_\infty(u) < 0$) and zero at $P_p$, for every irreducible
polynomial $p$ that divides $u$ (since $v_p(u) > 0$). Thus,
$(u)_\infty = \deg(u) P_\infty$ and $\deg((u)_\infty)=\deg(u)$.
Also, it turns out that $deg(P_p)=\deg(p)$. Thus, $\sum_p v_p(u)$ is the
number of irreducible factors $u$ has, and $\sum_p v_p(u) \deg(P_p)$
is the degree of $u$. In total, $\deg((u))=0$.
\end{example}

For a divisor $G$, we define the \emph{Riemann-Roch space} $\cL(G)$
to be:
$$\cL(G) = \set{ x\in F \;:\; (x) \ge -G} \cup \set{0}\,.$$


\begin{example}[Continued from Examples~\ref{exa:rational} and~\ref{exa:rational:2}]
%
%
%
For the divisor $G=kP_\infty$ the Riemann-Roch space $\cL(G)$ is the set of all polynomials of
degree at most $k$ over $\Fq$.
\end{example}

The set $\cL(G)$ is a vector space and furthermore has finite
dimension. We define the dimension of $G$ by $\dim (G) = \dim
\cL(G)$ and we use the two notations interchangeably. The fact that
the degree of each principal divisor is $0$ implies that if $\deg(G)
< 0$ then $\dim(\cL(G)) = 0$.

\subsubsection{Geometric Goppa codes}\label{subsec:GoppaCodes}
A Goppa code is specified by a triplet $(F,\Pl,G)$, where $F/\Fq$
is a function field, $\Pl = \set{P_1,\ldots,P_n}$ is a set of
places of degree $1$ and $G$ is an arbitrary divisor with no
support over any place in $\Pl$. Notice that for any $x \in
\cL(G)$, $v_{P_i}(x) \ge 0$ and therefore $x \in O_{P_i}$ and $x
(P_i) \in \Fq $. The triplet $(F,\Pl,G)$ specifies the code:
$$C(\Pl;G)=\set{(x(P_1),\ldots,x(P_n)) \;:\; x \in \cL(G)} \subseteq \Fq ^n\,.$$

\begin{claim}[{\cite[Cor II.2.3]{S:93}}]\label{clm-goppa-code}
If $\deg(G) < n$ then $C(\Pl;G)$ is an $[n,\dim(G),n-\deg(G)]$
linear code over $\Fq $.
\end{claim}


We call $n-\deg(G)$ the \emph{designated} distance of the code.
Given the designated distance $d$, we would like to find a code $G$
with that designated distance and maximal dimension, \ie, maximize
$\dim(G)$. It turns out that for every $G$, $\dim(G) \le \deg(G)+1$
(which is the Singleton bound in coding theory). Our goal is to
minimize the gap between $\dim(G)$ and $\deg(G)$. It turns out that
for any function field $F/\Fq$ there exists a constant $g \in
\mathbb{N}$, such that for any divisor $G \in \cD_F$, $\deg(G) -
\dim(G) \le g-1$. The minimal integer with this property is called
the \emph{genus} of $F/\Fq $. The Riemann-Roch Theorem states that:

\begin{theorem}[{\cite[Thm I.5.17]{S:93}}] If $\deg(G) \ge 2g-1$ then
$\dim(G)=\deg(G)-g+1$.
\end{theorem}

This, in particular, allows one to easily compute the dimension of
the code when $\deg(G) \ge 2g-1$. The only remaining question is
whether there are function fields $F/\Fq$ with a large number
$N=N(F)$ of evaluation points (\ie, degree $1$ places) and a small
genus $g$. A \emph{negative} answer to that question is given by the Hasse-Weil bound:
\begin{theorem}[{Hasse-Weil bound~\cite[Thm V.2.3]{S:93}}]
 \label{thm:hasse-weil}
Let $F/\Fq$ be a function field of genus $g$. Then, the number $N$ of places of degree one satisfies $N \le (q+1)+2\sqrt{q}g$.
\end{theorem}


The Drinfeld-Vl\u{a}du\c{t} bound tells us that when $g$ tends to
infinity, the bound can be strengthened by about a factor of $2$,
and roughly speaking, $N \le g (\sqrt{q}-1)$. This is tight for
prime power squares $q$.

On the positive side, the good news %
is that following much
research, there are several beautiful explicit constructions that
meet the Drinfeld-Vl\u{a}du\c{t} bound, and we refer the interested
reader to the beautiful survey paper~\cite[Chapter 1]{GS:06}.

In this paper we look at divisors $G$ whose degree is smaller than
the genus. Much less is known about such small-degree divisors. In
this regime, $\dim(G)$ depends on the divisor $G$ itself, and not
only on its degree, as is the case when $\deg(G) \ge 2g-1$. For some
special algebraic function fields the vector space $\cL(G)$ (and
therefore also its dimension) is known in full. We discuss this
below.


\subsection{Concatenating AG codes with Hadamard}
\label{sec:hadamard}
In this section we consider the concatenation of an outer code with the Hadamard code.
If the outer code is an $[n_1,k_1,d]_q$ code and $q$ is a power of two, then
concatenating it with the $[q,\log q,q/2]_2$ Hadamard code gives an
$[n=n_1 q, k=k_1 \log q]_2$ code that is $\eps={(n_1-d)}/{n_1}$
balanced, because non-zero symbols in the outer code expand by the
concatenation to perfectly balanced blocks.\footnote{If $q$ is a
power of $2$, then the resulting concatenated code is linear.
Concatenation is well defined even when $q$ is not a power of $2$.
In such a case we embed $\Fq$ into $\Ftwo^{\lceil \log q \rceil}$
using any one-to-one mapping. The resulting (non-linear) code has
essentially the same dimension and distance as in the previous case
-- the only difference is a small loss due to the fact that
$2^{\lceil \log q \rceil}$ is slightly larger than $q$. {}From now on
we will discuss the simpler case where $q$ is a power of two,
keeping in mind that everything also holds for arbitrary $q$.}
%

Using a $[q,k_1,q-k_1+1]_q$ Reed-Solomon code as the outer code, one
gets an $[n=q^2,k=k_1 \log q]_2$ code that is $\eps <{k_1}/{q}$
balanced. Rearranging the parameters, this gives an $[n,k]_2$
$\eps$-balanced code with
\begin{equation}\label{eq:aghp-params}
    n=O \left({k \over \eps \log (\frac{k}{\eps})} \right)^2\,.
\end{equation}
This is one of the constructions in~\cite{AGHP:92}.

Taking the outer code to be an $[N,\dim(G),N-\deg(G)]_q$ AG code
$C(\Pl;G)$ over $\Fq $ and concatenating it with Hadamard, we get an
$[n=Nq,k=\dim(G) \log q]_2$ code that is $\eps={\deg(G)}/{N}$
balanced. We can choose an AG code which uses a curve of genus $g$
with $N = \Theta(g \sqrt{q})$ degree $1$ places (the asymptotic is
over $g$ going to infinity). Picking the divisor $G$ to be of degree
$\deg(G) \ge 2g$ and setting $q= {1}/{\eps^2}$ results in
%
%
$$ N = \frac{\deg(G)}{\eps} = \frac{\dim(G)+g-1}{\eps}
= \Theta\left( \frac{k}{\eps \logeps}\right)\,,$$ where the second equality
follows from the Riemann-Roch Theorem. Thus, we get an
$\eps$-balanced code of length
\begin{equation}\label{eq:regular-ag-params}
    n = Nq = O\left(\frac{k}{\eps^3 \logeps}\right)\,.
\end{equation}

In fact, if one takes an AG code over $\Fq$ with large genus $g \ge
\sqrt{q}$ then $$N \ge \frac{\dim(G)}{\eps} = \frac{k}{\eps \log
q}\quad\text{and}\quad q=\Omega\left(\frac{1}{\eps^2}\right)$$ and
equation~(\ref{eq:regular-ag-params}) is tight. Taking an AG code
with a small genus $g \le \sqrt{q}$ is essentially equivalent to
taking a Reed-Solomon outer code and cannot be better (up to
constant factors) than equation~(\ref{eq:aghp-params}). In what
follows, we show one can improve on both bounds when the AG code has
degree much smaller than the genus.


So we now turn our attention to the case where $\deg(G) \le 2g-1$.
In this case $\dim(G)$ depends on the divisor $G$ and not just its
degree. One special case is the case where $G = rQ$, $r \in
\mathbb{N}$ and $Q$ is a place of degree 1. For any such $r$,
$\dim(rQ)$ is either equal to $\dim((r-1)Q)$ or to $\dim((r-1)Q)+1$.
In the former case $r$ is said to be a \emph{gap number} of $Q$. The
Weierstrass Gap Theorem~\cite[Thm I.6.7]{S:93} says that for any
place $Q$ there are exactly $g={\rm genus}(F/\Fq )$ gap numbers, and
they are all in the range $[1,2g-1]$.

The non-gap numbers (also called \emph{pole numbers}) form a
semigroup of $\mathbb{N}$ (\ie, a set that is closed under
addition). This semigroup is sometimes referred to as the
\emph{Weierstrass semigroup} of $Q$. We say that a semi-group $S$ is
\emph{generated} by a set of elements $\set{g_i}$, if each $g_i \in S$ and,
furthermore, every element $s \in S$ can be expressed as $s=\sum a_i
g_i$ with $a_i \in \mathbb{N}$.

The structure of the Weierstrass semigroup is crucial to our
construction. We know that there are exactly $g$ non-gap elements of
this semigroup in the range $[1,2g]$. If these elements are too
concentrated on the upper side of the range then the behavior of
$\dim(rQ)$ will be very similar to the case where $r > 2g-1$. Thus,
our goal is to find a function field $F$ that has many places of
degree $1$, say, $N(F) \ge \Omega(g \sqrt{q})$, while at the same
time $F$ has a degree $1$ place $Q$ with a ``good'' Weierstrass
semigroup.

\subsection{The construction}
Let $p$ be a prime power and $q=p^2$. The Hermitian function field
over $\Fq $ (see~\cite[Lemma VI.4.4]{S:93}) can be represented as
the extension field $\Fq (x,y)$ of the rational function field
$\Fq(x)$ with $y^p + y = x^{p+1}$. This function field has $1+p^3$
places of degree one. First, there is the common pole $Q_\infty$ of
$x$ and $y$. Moreover, for each pair $(\alpha, \beta)\in \Fq$ with
$\beta^p + \beta = \alpha^{p+1}$ there is a unique place
$P_{\alpha,\beta}$ of degree one such that $x(P_{\alpha,\beta}) =
\alpha$ and $y(P_{\alpha,\beta}) = \beta$ and we already saw there
are $p^3$ such pairs. The genus of the Hermitian function field is
$g=p(p-1)/2$.

For the outer code we take the Goppa code $C_r =
C(\Pl,G=rQ_\infty)$, where $\Pl$ is the set of all degree $1$ places
$P_{\alpha,\beta}$ mentioned above and $r=\eps p^3$.
%
%
%
%
The Weierstrass semigroup of $G$ is generated by $p$ and $p+1$,
and a basis for $\cL(G)=\cL(r Q_\infty)$ is $$\set{x^i y^j \;:\;
\text{$j\le p-1$ and $ip+j(p+1)\le r$}}\,.$$ The dimension of the
code is
$$\big|\{(i,j)\;:\; \text{$j\le p-1$ and $ip+j(p+1)\le r$}\} \big|\,.$$

We can now see the similarity between this construction and the one
in Section 2. The parameter $r$ is chosen such that
the constraint $ip+j(p+1)\le r$ forces  $j \le p-1$. Therefore, both
use evaluations of low degree bivariate polynomials over the same
set of $p^3$ points.\footnote{The only slight difference is that in
this construction we take all bivariate polynomials with bounded
\emph{weighted} total degree. However, the weight is nearly
identical for both variables and so this does not affect much the
parameters of the construction.}

\begin{theorem}
For every $k$ and every $\eps$ such that 
\[
\frac{\eps}
{\sqrt{\log(1/\eps)}} \le \frac{1}{\sqrt{k}}\,,
\]
there exists an
explicit $[n,\Omega(k)]_2$ code that is $\eps$-balanced, with
\[
n=O\left(\frac{k}{\eps^2 \log(1/\eps)}\right)^{5/4}\,.
\]
\end{theorem}

\begin{proof}
For a given $k$ and $\eps$, let $$p \in \left[\half
\left(\frac{k}{\eps^2 \log(1/\eps)}\right)^{1/4}, \;
\left(\frac{k}{\eps^2 \log(1/\eps)}\right)^{1/4}\right]$$ be a power
of two. It can verified that 
\begin{align*}
\frac{1}{16p^4}\; &\le\; \eps \;\le\;
\frac{1}{p}\\
\intertext{as}
  \frac{1}{p} \;&\;\ge \left(\frac{\eps^2 \log(1/\eps)}{k}\right)^{1/4} \;\ge\;
\left(\eps^2 \log(1/\eps) \cdot \frac{\eps^2}{\log(1/\eps)}\right)^{1/4}
\;=\; \eps\\
\intertext{and}
  \eps \;&=\; \eps^2 \cdot \frac{1}{\eps} \;\ge\; \eps^2 \cdot \log\left(\frac{1}{\epsilon}\right)
\;\ge\; \frac{k}{16 p^4} \;\ge\; \frac{1}{16 p^4}\,,
\end{align*}
and so $\log(1/\eps) = \Theta(\log(p))$.
%
%
%
%
%
%
%

Let $r=\eps p^3$ and let $\Fq $ be the field with $q=p^2$ elements.
Let $F$ denote the Hermitian function field over $\Fq $ and let
$\Pl$ denote its set of places of degree 1, excluding $Q_\infty$.
This implies that $|\Pl| = p^3$. Define the divisor $G$ to be $G =
rQ_\infty$. Since $r \le p^2$,
$$\dim (r Q_\infty) \ge \left(\frac{r}{2(p+1)}\right)^2 = \Omega\left(\eps^2
p^4\right)=\Omega\left(\frac{k}{\log(p)}\right)\,.$$ By
\expref{Claim}{clm-goppa-code}, the Goppa code that is obtained from the
triplet $(F, \Pl, G)$ is a
$$\left[p^3,\Omega\left(\frac{k}{\log(p)}\right),p^3 - r\right]_{p^2}$$ code. Concatenating this code with Hadamard gives a $[p^5,\Omega({k})]_2$ code that is $\eps$-balanced (since
${r}/{p^3}=\eps$). Now, by our choice of $p$, it follows that
$${k \over \eps^2 \logeps} = \Theta(p^4)$$ and therefore 
\[
n=p^5
=O\left(\left(\frac{k}{\eps^2 \logeps}\right)^{5/4}\right)
\]
as desired.
\end{proof}



\section{Limits of the approach}
\label{sec:voloch}

As explained in \expref{Section}{subsec:GoppaCodes}, the genus measures
the maximal loss in dimension compared to the degree. The
Drinfeld-Vl\u{a}du\c{t} bound implies that the number of evaluation
points (which is bounded by the number of degree one places $N(F)$)
is at most $O(g \sqrt{q})$ when $N(F) \gg q$. In
\expref{Section}{sec:hadamard} we saw this implies that when $\deg(G) >
2g$, concatenating the best AG code $C(Y;G)$ with the Hadamard code
cannot give $\eps$-balanced codes of dimension $k$ and length
\[
n=o\left(\frac{k}{\eps^3 \log(1/\eps)}\right)\,.
\]

Our construction shows that substantially better results are
possible when $\deg(G) \ll g$. Namely, we show that there exists a
code $C(Y;G)$ with $\deg(G) \ll g$ such that when this code is
concatenated with Hadamard, it gives a 
%
$k$-dimensional
$\eps$-balanced code of length 
\[
n=O\left(\left(\frac{k}{\eps^2
\log(1/\eps)}\right)^{5/4}\right)\,.
\]
It is therefore natural to ask what
are the limits of our approach. More concretely we ask what are the
best codes one can construct by concatenating an AG code with a
Hadamard code? Let us state the question precisely. We look at
constructions of the following structure:

\begin{itemize}
\item
An outer AG code $C=C(\Pl;G)$, defined by an algebraic function
field $F/\Fq$, a set of degree $1$ places $\Pl$ and a divisor $G \in
\cD_F$ with no support over any place in $\Pl$.

\item An inner Hadamard code.
\end{itemize}

In the analysis we view $C$ as a $[|\Pl|,\dim(G),|Y|-\deg(G)]_q$
code, and then the concatenated code has parameters
\[
\left[|Y|
q,\; \dim(G) \log(q),\; \half-\frac{\deg(G)}{|Y|}\right]_2\,.
\]
 Notice that it may
be the case that $C$ has better distance than the so-called
\emph{designated} distance, but as far as we are concerned the
analysis does not take advantage of that, and we take the distance
to be $|Y|-\deg(G)$.

In this section we prove:
\begin{theorem}
\label{thm:lower} Any $\eps$-balanced $[n,k]_2$ code that is
constructed and analyzed as above, must have
$$n \ge \Omega\left(\frac{k}{\eps^2} \cdot \min\set{\frac{k}{\logepsksquare},\frac{1}{\sqrt{\eps} \logepsk}}\right)\,.$$
\end{theorem}


For the proof we need definitions and theorems about finite
extensions of algebraic functions fields. Specifically, for an
extension $F$ of a function field $F'$, we use the following
notation:

\begin{itemize}
\item
A place $P \in \AllPl_{F}$ \emph{lying over} a place $P' \in \AllPl_{F'}$,
denoted by $P|P'$, see~\cite[Def III.1.3]{S:93},
\item
The \emph{ramification index} of $P$ over $P'$, denoted by $e(P|P')$, see~\cite[Def III.1.5]{S:93},
\item
The \emph{conorm} of a divisor $G' \in \cD_{F'}$, denoted by
$\con_{F/F'}(G')$, see~\cite[Def III.1.8]{S:93}.
\end{itemize}

For more details we refer the reader to~\cite[Chapter III]{S:93}.

\subsection{AG theorems about degree vs.\ dimension}

It turns out that the above question boils down to the question of
whether there are function fields with many degree $1$ places
(compared to the genus) and with low-degree divisors (of degree much
smaller than the genus) of high dimension. We start by presenting
two AG theorems relating degree to dimension in the small degree
regime (when the degree is smaller than the genus).

The first argument we present shows that any divisor with
non-trivial dimension must have degree at least $N(F)/(q+1)$. The
argument was shown to us by Henning Stichtenoth~\cite{S:09}.

\begin{lemma}
\label{lem:henning} Let $F/\Fq $ be a function field and $G \in
\cD_F$ a divisor with $\dim(G) > 1$. Then $N(F) \le \deg(G) \cdot
(q+1)$.
\end{lemma}

\begin{proof}
As $\dim(G) > 1$, there exists some $x \in F \setminus \Fq$ such
that $(x) \ge -G$. Fix any such $x$. In particular,
$\deg((x)_\infty) \le \deg(G)$. Also, by~\cite[Thm I.4.11]{S:93},
$\deg(x)_\infty=[F:\Fq (x)]$. We may view $F$ as a finite extension
over the rational function field $\Fq (x)$. Every place of degree
$1$ of $F$ lies above some place of degree $1$ of $\Fq (x)$. There
are exactly $q+1$ places of degree $1$ of $\Fq (x)$, and each one of
them may split to at most $[F:\Fq (x)]$ places of degree $1$ of $F$
(by the fundamental equality,~\cite[Thm III.1.11]{S:93}).
Altogether, $N(F) \le (q+1) [F:\Fq (x)] = (q+1) \deg(x)_\infty \le
(q+1) \deg(G)$.
\end{proof}

\begin{remark}
\label{rem:henning} \expref{Lemma}{lem:henning} only uses the fact that
$G$ is non-trivial. We wonder if one can strengthen the lemma for
divisors $G$ of high dimension. In particular, is it true that if
$\dim(G) > \ell$ then
\[
N(F) \le \frac{\deg(G) \cdot (q+1)}{f(\ell)}
\]
for some function $f$ that goes to infinity with $\ell$?
\end{remark}

%
%
%
%
%

We now move to the second theorem. For a set $S \subseteq F$, where
$F$ is a field, let $\closure(S)$ denote the minimal subfield of $F$
that contains $S$. Following our work, Voloch~\cite{V:09} showed,
based on the Castelnuovo bound, that:
\begin{theorem}[{\cite[based on the Castelnuovo bound]{V:09}}]
\label{thm:castel} Let $K$ be an arbitrary field. Let $F/K$ be a
function field of genus $g$. Let $G \in\cD_F$ be a divisor with
degree $d+1$ and dimension $\ell+2$ such that
$\closure(\cL(G))=F$. Let $m = d \;\div\; \ell$ and $r = d \bmod
\ell$. Then $$g \le m(m-1)\ell + m(2r+1)\,,$$ and, in particular, $g
\le m(m+1)\ell$.
\end{theorem}

Using \expref{Theorem}{thm:castel} requires an assumption on the AG
code, namely, that the closure of the Riemann space of the divisor
used to define the code is the entire function field $F$. The
following lemma, based on private communication with Voloch, shows
that this assumption is inessential when analyzing the rate versus
distance problem.
\begin{lemma}
\label{lem:closure-assumption} Let $K$ be a finite field, $F/K$ be a
function field, $G \in\cD_F$ is a divisor. Let $C$ be the Goppa code
of length $n$, dimension $k$, and designated relative distance
$\delta$, specified by some triplet $(F,Y,G)$. Define a function
field $F' = \closure(\cL(G))$. Then there exists a Goppa code $C'$
defined by a triplet $(F',Y',G')$, of length $n' \le n$, dimension
$k$ and designated relative distance $\delta' \ge \delta$, such that
$\closure(\cL(G'))=F'$.
\end{lemma}
\begin{proof}
%
%
%
%




We first define the new triplet $(F',Y',G')$.

\begin{itemize}
\item
We already have that $F' = \closure(\cL(G))$.

\item
Next let $B=\set{s_1,\ldots,s_{k}}$ be a basis for $\cL(G)$. Define,
$$G'' = \sum_{P' \in \AllPl_{F'}} \max_i
\set {-v_{P'}(s_i)} \cdot P'\,.$$ We would like to exchange $G''$ with
an equivalent divisor that has no support over places of degree 1.
By the Weak Approximation Theorem~\cite[Thm I.3.1]{S:93} there
exists $z\in F/K$ such that for every place $P$ of degree $1$,
$v_P(z) = -v_P(G)$ and we let $$G' = G'' + (z)\,.$$

\item
Define a set $Y' \subset \AllPl_{F'}$ by
$$ Y' = \set{ P' \in \AllPl_{F'} \;:\; \text{$\exists P \in Y$ such
    that $P\mid P'$}}\,.$$
Observe that since $Y$ consists only of places of degree 1 this is
also true for $Y'$ (see~\cite[Proposition III.1.6]{S:93}).
\end{itemize}

Consider the Goppa code $C'$ defined by the triplet $(F',Y',G')$.
Notice that $Y'$ does not intersect $G'$ because $Y'$ contains only
degree $1$ places and $G'$ has no support over degree $1$ places. We
will prove:
\begin{itemize}
  \item The dimension of $C'$ is the same as $C$, \ie, $\dim(\cL_F'(G'))=k$.
  \item The length of $C'$ is at most the length of $C$, \ie, $n'=|Y'| \le |Y| = n$.
  \item The designated relative distance of $C'$ is at least as good as in $C$, \ie,
$$\delta' = 1 - \frac{\deg(G')}{|Y'|} \ge \delta\,.$$
\end{itemize}

For the proof we will show:

\begin{claim}
\label{cl:conorm}
$$G \ge \con_{F/F'}(G'')\,.$$
\end{claim}

With that we can prove the three assertions above about $C'$:

\begin{description}
\item [Dimension:]
Since $\cL_{F'}(G'') \subseteq \cL_F(\con_{F/F'}(G''))$, it follows
that $$\dim(\cL_{F'}(G'')) \le \dim(\cL_F(\con_{F/F'}(G'')))\,.$$ Thus,
by \expref{Claim}{cl:conorm},
$$\dim(\cL_F(G)) \ge \dim( \cL_F(\con_{F/F'}(G''))) \ge \dim(\cL_{F'}(G'')) \ge
|B| = \dim(\cL_F(G))\,,$$ and therefore
$$\dim(\cL_{F'}(G''))=\dim(\cL_F(G))=k\,.$$
The claim follows since $\dim(\cL_{F'}(G'))=\dim(\cL_{F'}(G''))$ by~\cite[Lemma I.4.6]{S:93}.

\item [Length:] $|Y| \ge |Y'|$, since every place of $Y$ lies over exactly one place of $Y'$,
see~\cite[Proposition III.1.7]{S:93}.

\item [Designated distance:]
By \expref{Claim}{cl:conorm}, $\deg(G) \ge \deg( \con_{F/F'}(G''))$. By~\cite[Cor III.1.13]{S:93},
$$\deg(\con_{F/F'}(G''))=[F:F'] \cdot \deg(G'')\,.$$ Since
$\deg(G'')=\deg(G')$ it follows that $$\deg(G) \ge [F:F'] \cdot
\deg(G')\,.$$

Also, since every place $P'$ can split to at most $[F:F']$ places in
$F/K$ we have $$|Y| \le |Y'| \cdot [F:F']\,.$$

Altogether,
$$\delta' = 1 - \frac{\deg(G')}{|Y'|}
\ge  1 - \frac{\deg(G)}{[F:F'] \cdot |Y'|} \ge 1 -
\frac{\deg(G)}{|Y|} = \delta\,.$$
\end{description}
\end{proof}

We are left with proving \expref{Claim}{cl:conorm}.

\begin{proof}[Proof of \expref{Claim}{cl:conorm}]
 By definition, $B \subseteq \cL(G'')$, and therefore
$$F' = \closure(\cL(G)) = \closure(B) \subseteq \closure(\cL(G''))\,,$$
and $\closure(\cL(G''))=F'$.

Also, for any $P|P'$ (where  $P' \in \cP_{F'}$ and $P \in \cP_F$)
and for any $i$, $$e(P|P') \cdot v_{P'}(s_i) = v_{P}(s_i) \ge
-v_{P}(G)\,,$$ where the last inequality is simply because $s_i \in
\cL(G)$. Therefore
$$v_{P'}(G'') = \max \set{-v_{P'}(s_i)} =
\max \set{- \frac{v_{P}(s_i)}{e(P|P')}} \le
\frac{v_{P}(G)}{e(P|P')}\,,$$ and the claim follows from the definition
of the conorm.
\end{proof}

\subsection{The bound}

We are now ready to prove \expref{Theorem}{thm:lower}.
\begin{proof}[Proof of \expref{Theorem}{thm:lower}]
Assume a code is obtained by concatenating the AG code specified by
the triplet $(F/\Fq,Y,G)$ with the Hadamard code. Let $\ell =
\dim(G)$ and $d = \deg(G)$. The AG code $C(Y;G)$ is a $[|Y|,\ell]_q$
code, with designated distance $|Y|-d$. The concatenated code is
therefore a $$[n=|Y| \cdot q, k=\ell \log(q)]_2$$ code which is
$\eps$-balanced for
$$\eps=\frac{d}{|Y|}\,.$$
By \expref{Lemma}{lem:closure-assumption} we can assume without loss of
generality that $\closure(\cL(G))=F$.


There are two extreme cases that we handle separately:

\begin{description}
\item [Large base field:]
If the base field size $q$ is too large the theorem is trivially
true. Namely, If $q>\frac{k}{\eps^3}$ we are done because
$n \ge q > {k}/{\eps^3}$. We can therefore assume without loss
of generality that
\begin{align}
\nonumber q \;& \le \; \frac{k}{\eps^3}\\
\intertext{and}
\label{eq:qlogeps}\log(q)\; & = \;O\left(\log\left(\frac{k}{\epsilon}\right)\right)\,.
\end{align}

\item [Few evaluation points:]
If the number of evaluation points is about the field size, we are
essentially in the Reed-Solomon case and we are done. Specifically,
if $|Y| \le 4q$ then $$4n= |Y| \cdot 4q \ge |Y|^2 =
\frac{d^2}{\eps^2} \ge \frac{ \ell^2 }{\eps^2} = \frac{ k^2 }{\eps^2
\log^2(q)} = \Omega\left(\frac{k^2}{\eps^2 \logepsksquare}\right)\,,$$
and we are done. We can therefore assume without loss of generality
that
\begin{eqnarray*}
|Y| > 4q\,.
\end{eqnarray*}
This also implies that $\sqrt{q} < g$ since by
\expref{Theorem}{thm:hasse-weil}, $|Y|  \le  N(F) \le  q+1 +2g
\sqrt{q}$. We can therefore conclude (again, by
\expref{Theorem}{thm:hasse-weil}) that
\begin{eqnarray}\label{eq:N}
N(F) \le 4g\sqrt{q}\,.
\end{eqnarray}
\end{description}






Let $m = d \;\div\; \ell \ge 1$. By \expref{Theorem}{thm:castel}, $g \le
2m^2 \ell$, and by equation (\ref{eq:N}), $N(F) \le 8m^2 \ell \sqrt{q}$.
Thus,
\begin{equation*}
n \;=\; |Y| \cdot q \;=\; \frac{N(F) \cdot |Y| \cdot q}{N(F)} \;\ge\;
\frac{N(F) \cdot |Y| \cdot q}{8m^2 \ell \sqrt{q}}\,.
\end{equation*}
Substituting $m \le d/\ell$ and $d=\eps |Y|$ we see that
\begin{equation*}
n \;\ge\;  \frac{N(F) \cdot \sqrt{q} \cdot \ell}{8 \eps^2 |Y|}.
\end{equation*}
Substituting $\ell = {k}/{\log(q)}$ and using $N(F) > |Y|$,
\begin{equation*}
n\;\ge\;  \frac{\sqrt{q} \cdot k}{8 \eps^2 \log (q)} \;=\;
\Omega\left(\frac{\sqrt{q} \cdot k}{\eps^2 \logepsk}\right),
\end{equation*}
where the last equality follows from equation~(\ref{eq:qlogeps}). To
finish the argument notice that by \expref{Lemma}{lem:henning}, $N(F)
\le d(q+1)$. This implies
\[
\frac{d}{\eps} =|Y| \le d(q+1)\quad \text{and} \quad \eps
\ge \frac{1}{q+1}\,,\]
hence,
\[
n =\Omega\left( \frac{k}{\eps^{2.5}
\logepsk}\right)\,.\qedhere
\]
\end{proof}

\begingroup\raggedright
%
%
%
%
%
\providecommand{\bibhead}[1]{}
%
\expandafter\ifx\csname pdfbookmark\endcsname\relax%
  \providecommand{\tocrefpdfbookmark}{}
\else\providecommand{\tocrefpdfbookmark}{%
   \hypertarget{tocreferences}{}%
   \pdfbookmark[1]{References}{tocreferences}}%
\fi

\tocrefpdfbookmark
\begin{thebibliography}{10}

\bibitem{ABNNR:92}\bibhead{ABNNR:92}
{\sc Noga Alon, Jehoshua Bruck, Joseph Naor, Moni Naor, and Ron~M. Roth}:
  Construction of asymptotically good low-rate error-correcting codes through
  pseudo-random graphs.
\newblock {\em IEEE Trans. Inform. Theory}, 38(2):509--516, 1992.
\newblock Preliminary version in
  \href{http://dx.doi.org/10.1109/ISIT.1991.695195}{ISIT'91}.
\newblock [\epfmtdoi{10.1109/18.119713}]

\bibitem{AGHP:92}\bibhead{AGHP:92}
{\sc Noga Alon, Oded Goldreich, Johan H{\aa}stad, and Ren\'{e} Peralta}: Simple
  constructions of almost $k$-wise independent random variables.
\newblock {\em Random Structures {\&} Algorithms}, 3(3):289--304, 1992.
\newblock Preliminary version in
  \href{http://dx.doi.org/10.1109/FSCS.1990.89575}{FOCS'90}.
\newblock [\epfmtdoi{10.1002/rsa.3240030308}]

\bibitem{focsversion}\bibhead{focsversion}
{\sc Avraham Ben-Aroya and Amnon Ta-Shma}: Constructing small-bias sets from
  algebraic-geometric codes.
\newblock In {\em Proc. 50th FOCS}, pp. 191--197. IEEE Comp. Soc. Press, 2009.
\newblock [\epfmtdoi{10.1109/FOCS.2009.44}]

\bibitem{F:08}\bibhead{F:08}
{\sc William Fulton}: {\em Algebraic Curves: An Introduction to Algebraic
  Geometry}.
\newblock Third edition, 2008.
\newblock \href{http://www.math.lsa.umich.edu/~wfulton/CurveBook.pdf}{Author's
  version}.

\bibitem{GS:06}\bibhead{GS:06}
{\sc Arnaldo Garcia and Henning Stichtenoth}, editors.
\newblock {\em Topics in Geometry, Coding Theory and Cryptography}.
\newblock Volume~6.
\newblock Springer, 2007.
\newblock Available from
  \href{http://www.springer.com/mathematics/numbers/book/978-1-4020-5333-7}{Springer}.

\bibitem{L:02}\bibhead{L:02}
{\sc Serge Lang}: {\em Algebra}.
\newblock Springer, revised third edition, 2002.
\newblock Available from
  \href{http://www.springer.com/mathematics/algebra/book/978-0-387-95385-4}{Springer}.

\bibitem{NN:93}\bibhead{NN:93}
{\sc Joseph Naor and Moni Naor}: Small-bias probability spaces: Efficient
  constructions and applications.
\newblock {\em SIAM J. Comput.}, 22(4):838--856, 1993.
\newblock Preliminary version in
  \href{http://dx.doi.org/10.1145/100216.100244}{STOC'90}.
\newblock [\epfmtdoi{10.1137/0222053}]

\bibitem{S:93}\bibhead{S:93}
{\sc Henning Stichtenoth}: {\em {Algebraic Function Fields and Codes}}.
\newblock Springer, 1993.
\newblock Available from
  \href{http://www.springer.com/mathematics/algebra/book/978-3-540-76877-7}{Springer}.

\bibitem{S:09}\bibhead{S:09}
{\sc Henning Stichtenoth}: Private communication, 2009.

\bibitem{TVZ:82}\bibhead{TVZ:82}
{\sc Michael~A. Tsfasman, Serge~G. Vl{\u{a}}du{\c{t}}, and Thomas Zink}:
  {Modular curves, {S}himura curves, and {G}oppa codes, better than
  {V}arshamov-{G}ilbert bound}.
\newblock {\em Mathematische Nachrichten}, 109(1):21--28, 1982.
\newblock [\epfmtdoi{10.1002/mana.19821090103}]

\bibitem{V:09}\bibhead{V:09}
{\sc Jos\'{e}~Felipe Voloch}: Special divisors of large dimension on curves
  with many points over finite fields.
\newblock {\em Portugaliae Mathematica}, 68(1):103--107, 2011.
\newblock [\epfmtdoi{10.4171/PM/1882}]

\end{thebibliography}
\endgroup
\end{document}
